% ==================================================================
% id:       seq-027
% topic:    סדרות
% subtopic: גבולות במובן הרחב
% solution: full
% ==================================================================
\textbf{חשבו את גבול הסדרות הבאות ונמקו את צעדיכם.}
\begin{enumerate}\bMs{2}
\item $\ds\lim_{n \to \infty} \frac{\lr{2 - (-1)^n}\lr{n^3 + n}}{1 + 2n^2}$
\hint{הגורם $2-(-1)^n$ חסום מלרע על ידי מספר חיובי. השתמשו בזה כדי להשוות לסדרה פשוטה יותר.}
\begin{solbox}
    הגורם הראשון במכפלה שבמונה הוא סדרה מתחלפת בין המספרים
    $1,3$,
    לכן בכל מקרה
    $$1\le 2-(-1)^n\le 3$$
    ולכן
    $$\frac{3(n^3+n)}{1+2n^2}\le\frac{\lr{2 - (-1)^n}\lr{n^3 + n}}{1 + 2n^2}\le
    \frac{3(n^3+n)}{1+2n^2}$$
    כאשר נשים לב ששתי הסדרות שקיבלנו מימין ומשמאל הן שואפות ל-$\infty$, לכן מהאי-שיוויון השמאלי (ללא צורך בימני) נובע מכלל הפיצה כי הסדרה מתכנסת ל-$\infty$.
\end{solbox}

%%%%%
\item $\ds\lim_{n \to \infty} \frac{5n + \sqrt[6]{n^7 - n^5 - n + 1}}{\sqrt[4]{1 + 2n + n^4} + \sqrt[3]{n^2 + 5}}$
\hint{חלקו את המונה והמכנה בחזקה הגבוהה ביותר של $n$ שמופיעה במכנה (שימו לב לחזקות השבריות שבתוך השורשים).}
\begin{solbox}
נסמן את איברי הסדרה ב-
$a_n$.
מתקיים
\begin{align*}
a_n
=\ \frac{5n + n^{\frac{7}{6}}\sqrt[6]{1 - \frac{1}{n^2} - \frac{1}{n^6} + \frac{1}{n^7}}}{n \sqrt[4]{\frac{1}{n^4} + \frac{2}{n^3} + 1} + n^{\frac{2}{3}}\sqrt[3]{1 + \frac{5}{n^2}}}
=\ \frac{5 + n^{\frac{1}{6}}\sqrt[6]{1 - \frac{1}{n^2} - \frac{1}{n^6} + \frac{1}{n^7}}}{\sqrt[4]{\frac{1}{n^4} + \frac{2}{n^3} + 1} + \frac{1}{n^{\frac{1}{3}}}\sqrt[3]{1 + \frac{5}{n^2}}}.
\end{align*}
לאחר שחילקנו את המונה והמכנה בחזקה הגדולה ביותר של
$n$
מהמכנה,
נשים לב כי הגבול של המונה הוא
$\infty$
והגבול של המכנה הוא
$1$,
לכן
לפי אריתמטיקה של גבולות במובן הרחב, נסיק כי
$a_n\goes\infty$.
\end{solbox}
%%%%%
\item $\ds\lim_{n \to \infty} n + \sqrt{n^2 - n}$
\begin{solbox}
מתקיים
$$,n + \sqrt{n^2 - n}\geq n \goes\infty$$
מכאן, לפי כלל הפיצה, שגבולה של הסדרה הוא אינסוף.
\end{solbox}
%%%%%
\eMs
\end{enumerate}
